% P-DINA secara mendalam: model diagnosis kognitif yang diimplementasikan di proyek CDA-CDM.
% Build (dari folder slides/): pdflatex pdina_deep.tex (dua kali)
\documentclass[11pt,a4paper]{article}
\usepackage[T1]{fontenc}
\usepackage[utf8]{inputenc}
\usepackage{lmodern}
\usepackage[indonesian]{babel}
\usepackage[margin=2.3cm]{geometry}
\usepackage{microtype}
\usepackage{amsmath,amssymb}
\usepackage{booktabs,tabularx,array,colortbl}
\usepackage{xcolor}
\usepackage{tikz}
\usetikzlibrary{arrows.meta,positioning,shapes.geometric,calc,fit,backgrounds}
\usepackage{pgfplots}
\usepackage[most]{tcolorbox}
\usepackage{caption}
\usepackage{enumitem}
\usepackage{titlesec}
\usepackage{needspace}
\usepackage{fancyhdr}
\usepackage[hidelinks]{hyperref}
% Shared colours, TikZ styles and pgfplots defaults (used by slides.tex and naskah.tex).
\definecolor{navy}{HTML}{1F3A5F}
\definecolor{accent}{HTML}{D9822B}
\definecolor{tealc}{HTML}{2A9D8F}
\definecolor{gapred}{HTML}{C8553D}
\definecolor{okgreen}{HTML}{3A8D5C}
\definecolor{midgray}{HTML}{8A94A6}
\definecolor{lightbg}{HTML}{EEF2F7}
\definecolor{insuf}{HTML}{C9CED8}

\pgfplotsset{
  compat=1.18,
  every axis/.append style={
    font=\scriptsize,
    axis line style={midgray},
    tick style={midgray},
    label style={font=\scriptsize},
    tick label style={font=\scriptsize},
    legend style={font=\scriptsize, draw=none, fill=none},
    grid style={midgray!25},
    every node near coord/.append style={text=black!80},
  },
}
% Indonesian number format everywhere (0,5 and 1.000)
\pgfkeys{/pgf/number format/.cd, use comma, 1000 sep={.}}

% ragged-right columns (avoid stretched inter-word spacing in narrow cells)
\newcolumntype{P}[1]{>{\raggedright\arraybackslash}p{#1}}
\newcolumntype{L}{>{\raggedright\arraybackslash}X}

\tikzset{
  box/.style={draw=navy, rounded corners=3pt, fill=white, align=center, font=\scriptsize,
              inner sep=4pt, line width=0.6pt},
  boxnavy/.style={box, fill=navy, text=white},
  boxacc/.style={box, draw=accent, fill=accent, text=white},
  boxteal/.style={box, draw=tealc, fill=tealc, text=white},
  boxsoft/.style={box, draw=midgray, fill=lightbg},
  arr/.style={-{Stealth[length=2.2mm]}, line width=0.8pt, draw=navy!80},
  lbl/.style={font=\tiny, text=black!70, align=center},
}

% height used by every pgfplots figure; set before \input{figs/...}
\newlength\figH
\setlength\figH{5cm}


\captionsetup{font=small, labelfont={bf,color=navy}}
\titleformat{\section}{\Needspace{10\baselineskip}\large\bfseries\color{navy}}{\thesection}{0.6em}{}
\titleformat{\subsection}{\Needspace{7\baselineskip}\normalsize\bfseries\color{navy}}{\thesubsection}{0.6em}{}
\titlespacing*{\section}{0pt}{12pt}{4pt}
\titlespacing*{\subsection}{0pt}{8pt}{3pt}
\setlist{itemsep=1pt, topsep=3pt}
\pagestyle{fancy}
\fancyhf{}
\fancyhead[L]{\small\color{midgray}P-DINA secara mendalam}
\fancyhead[R]{\small\color{midgray}\thepage}
\renewcommand{\headrulewidth}{0.3pt}
\setlength{\parskip}{3pt}
\setlength{\parindent}{0pt}

\newtcolorbox{intinya}{colback=lightbg, colframe=navy, boxrule=0.4pt, arc=2pt, left=6pt, right=6pt,
  top=3pt, bottom=3pt, fontupper=\small}
\newtcolorbox{awas}{colback=accent!8, colframe=accent, boxrule=0.4pt, arc=2pt, left=6pt, right=6pt,
  top=3pt, bottom=3pt, fontupper=\small}
\newcommand{\sumber}[1]{{\footnotesize\color{midgray}#1}}

\begin{document}

\begin{center}
  {\LARGE\bfseries\color{navy} P-DINA secara Mendalam}\\[4pt]
  {\large Cara model diagnosis kognitif di proyek ini bekerja, langkah demi langkah}\\[6pt]
  {\small\color{midgray} Muhammad Fariq Faqih \quad·\quad \today}
\end{center}

\begin{intinya}
  \textbf{Tentang dokumen ini.} Semua angka berasal dari checkpoint yang dipakai aplikasi (\texttt{artifacts/model/})
  atau dihitung ulang dari checkpoint itu dengan kode proyek (\texttt{cdacdm/models.py}). Model juga dilatih ulang
  sekali untuk mencatat kurva ELBO; hasilnya identik dengan checkpoint (selisih $g$ dan $s$ = 0, 44 iterasi,
  32 detik CPU). Data: EEDI NeurIPS 2020 Task 3/4, yang merupakan data pengembangan (mayoritas siswa SMP di Inggris),
  bukan data kelas SMA penelitian.
\end{intinya}

% ======================================================================================================
\section{Pertanyaan yang dijawab P-DINA}

Seorang guru melihat lembar jawaban siswa: 13 dari 25 soal \emph{Writing Expressions} benar. Apakah siswa ini sudah
menguasai konsep itu? ``52\% benar'' tidak langsung menjawab, karena:
\begin{itemize}
  \item soal pilihan ganda bisa dijawab benar dengan \textbf{menebak};
  \item siswa yang paham pun bisa \textbf{salah karena ceroboh};
  \item setiap soal punya tingkat tebakan dan kecerobohan yang berbeda;
  \item soal yang tidak dikerjakan bukan berarti salah.
\end{itemize}
P-DINA menjawab pertanyaan itu dengan angka yang jelas artinya: \textbf{P(mastery)}, yaitu peluang siswa sudah
menguasai sebuah konsep, \emph{setelah} memperhitungkan semua hal di atas. Untuk siswa tadi jawabannya 0,34 (dihitung
lengkap di Bagian~\ref{sec:worked}).

P-DINA termasuk model \textbf{unsupervised} (model variabel laten): tidak ada data yang memberi tahu ``siswa ini
sebenarnya menguasai konsep ini''. Model hanya melihat jawaban benar/salah dan Q-matrix dari pakar, lalu mencari
pola penguasaan tersembunyi yang paling masuk akal untuk menjelaskan jawaban-jawaban itu.

% ======================================================================================================
\section{Bahan baku: jawaban siswa dan Q-matrix}

\textbf{Matriks jawaban.} Data train berisi 3.934 siswa $\times$ 948 soal. Hanya 29,8\% sel yang terisi (median
241 jawaban per siswa), sisanya kosong karena siswa tidak mengerjakan soal itu. Nilai sel: 1 = benar, 0 = salah,
kosong = tidak teramati. Total 1.112.450 jawaban teramati.

\textbf{Q-matrix.} Tabel biner 948 soal $\times$ 57 konsep yang disusun pakar: $q_{jk}=1$ bila soal $j$ menguji konsep
$k$. Q-matrix \emph{tidak dipelajari} model; ia adalah pengetahuan awal yang menentukan struktur model. Sangat jarang
terisi (1,78\%): 934 soal menguji tepat satu konsep, hanya 14 soal menguji dua konsep.

\begin{table}[h]
  \centering\small
  \begin{tabular}{@{}lr@{\hspace{2em}}lr@{}}
    \toprule
    Konsep target & Soal & Konsep target & Soal \\
    \midrule
    Writing Expressions & 25 & Factorising into a Single Bracket & 8 \\
    Naming Co-ordinates in 2D & 24 & Solving Linear Inequalities & 8 \\
    Substitution into Formula & 22 & Horizontal and Vertical Lines & 8 \\
    Simplifying Expr.\ (Collecting Like Terms) & 19 & Linear Equations & 7 \\
    Linear Sequences (nth term) & 12 & \textcolor{gapred}{Gradient Between Two Co-ordinates} & \textcolor{gapred}{2} \\
    Midpoint Between Two Co-ordinates & 9 & & \\
    \bottomrule
  \end{tabular}
  \caption{Jumlah soal per konsep target di Q-matrix. Konsep dengan 2 soal tidak pernah bisa memenuhi syarat bukti
  minimum 5 jawaban, sehingga tidak dapat didiagnosis.}
\end{table}

% ======================================================================================================
\section{Inti model: aturan DINA}

DINA adalah singkatan \emph{Deterministic Inputs, Noisy ``And'' gate}. Idenya dua langkah (Gambar~\ref{fig:tree}):

\begin{enumerate}
  \item \textbf{Apakah siswa ``siap'' untuk soal ini?} Siswa siap ($\eta_{ij}=1$) hanya jika menguasai \emph{semua}
        konsep yang dibutuhkan soal (gerbang ``AND''). Kurang satu konsep saja berarti belum siap.
        \[\eta_{ij}=\prod_{k}\alpha_{ik}^{\,q_{jk}}\]
        dengan $\alpha_{ik}\in\{0,1\}$ status penguasaan siswa $i$ pada konsep $k$.
  \item \textbf{Tambahkan ``noise'' yang realistis.} Siswa siap tetap bisa salah dengan peluang $s_j$ (\emph{slip});
        siswa belum siap tetap bisa benar dengan peluang $g_j$ (\emph{guess}):
        \[P(X_{ij}=1\mid\boldsymbol\alpha_i)=(1-s_j)^{\eta_{ij}}\,g_j^{\,1-\eta_{ij}}.\]
\end{enumerate}

\begin{figure}[h]
  \centering
  \resizebox{0.95\linewidth}{!}{% Illustration: how DINA turns a mastery profile + Q-matrix row into P(correct).
\begin{tikzpicture}[node distance=4mm and 7mm]
  \node[boxsoft, text width=3.1cm, align=left] (prof) {\textbf{Profil siswa} $\boldsymbol\alpha_i$\\[2pt]
     {\tiny Konsep A: menguasai (1)\\ Konsep B: belum (0)\\ Konsep C: menguasai (1)}};
  \node[boxsoft, text width=3.1cm, align=left, below=of prof] (item) {\textbf{Baris Q-matrix soal} $j$\\[2pt]
     {\tiny butuh A dan C (B tidak dipakai)}};
  \node[boxnavy, text width=2.7cm, anchor=west] (eta) at ($(prof.east)!0.5!(item.east)+(10mm,0)$)
     {$\eta_{ij}=\prod_k\alpha_{ik}^{q_{jk}}$\\[2pt]{\tiny $=1^1\cdot0^0\cdot1^1=1$ (``siap'')}};
  \draw[arr] (prof.east) -- (eta.west);
  \draw[arr] (item.east) -- (eta.west);
  \node[box, draw=okgreen, text width=3.6cm, anchor=west] (siap) at ($(eta.east)+(12mm,9mm)$)
     {\textbf{$\eta=1$ (siap)}\\[1pt]{\tiny P(benar) $=1-s_j$; salah hanya karena \emph{slip} (ceroboh)}};
  \node[box, draw=gapred, text width=3.6cm, anchor=west] (belum) at ($(eta.east)+(12mm,-9mm)$)
     {\textbf{$\eta=0$ (belum siap)}\\[1pt]{\tiny P(benar) $=g_j$; benar hanya karena \emph{guess} (menebak)}};
  \draw[arr] (eta.east) -- (siap.west);
  \draw[arr] (eta.east) -- (belum.west);
\end{tikzpicture}
}
  \caption{Cara DINA menghitung peluang benar. Soal butuh konsep A dan C; siswa menguasai keduanya, jadi siap.}
  \label{fig:tree}
\end{figure}

Setiap soal hanya punya dua parameter, $g_j$ dan $s_j$, dan keduanya punya arti langsung. Dengan nilai median di
EEDI ($g\approx0{,}32$, $s\approx0{,}22$), siswa yang menguasai menjawab benar sekitar 78\% dan yang belum sekitar 32\%.
Selisih keduanya, $1-s_j-g_j$, adalah \textbf{daya beda} soal: makin besar, makin jelas soal itu memisahkan siswa
yang paham dan yang belum. Rata-rata daya beda soal EEDI 0,41, tergolong sedang.

\subsection{Seperti apa soal-soal EEDI?}

Gambar~\ref{fig:gs} menempatkan ke-948 soal menurut guess dan slip hasil training. Soal di pojok kiri bawah adalah
soal ``tajam'' (jarang ditebak, jarang membuat siswa paham tergelincir); soal di dekat garis merah hampir tidak
membedakan siapa pun.

\begin{figure}[h]
  \centering
  \setlength\figH{6.2cm}
  \begin{minipage}[c]{0.52\linewidth}% Guess vs slip for all 948 items (slides/data/pdina_deep/item_gs.csv, from artifacts/model/item_parameters.csv)
\begin{tikzpicture}
\begin{axis}[
  width=\linewidth, height=\figH, xmin=0, xmax=0.72, ymin=0, ymax=0.72,
  xlabel={guess $g$ (peluang benar walau belum menguasai)}, ylabel={slip $s$ (peluang salah walau menguasai)},
  grid=major,
  xticklabel style={/pgf/number format/fixed, /pgf/number format/precision=1},
  yticklabel style={/pgf/number format/fixed, /pgf/number format/precision=1},
]
  \fill[gapred!10] (axis cs:0.9,0) -- (axis cs:0.72,0) -- (axis cs:0.72,0.72) -- (axis cs:0.18,0.72) -- cycle;
  \addplot[only marks, mark=*, mark size=0.7pt, navy, fill opacity=0.35, draw opacity=0.35]
    table[col sep=comma, x=guess, y=slip] {data/pdina_deep/item_gs.csv};
  \addplot[gapred, dashed, domain=0.18:0.72] {0.9-x};
  \addplot[okgreen, dashed, domain=0:0.4] {0.4-x};
  \node[font=\tiny, text=gapred, anchor=south west, align=left] at (axis cs:0.30,0.63) {daya beda $<0{,}1$\\(30 soal, nyaris tak berguna)};
  \node[font=\tiny, text=okgreen!70!black, anchor=south west, align=left, fill=white, inner sep=1pt] at (axis cs:0.005,0.005) {daya beda $>0{,}6$\\(81 soal, sangat tajam)};
\end{axis}
\end{tikzpicture}
\end{minipage}\hfill
  \begin{minipage}[c]{0.44\linewidth}\small
    \begin{itemize}
      \item Mayoritas soal berkumpul di $g\approx0{,}2$--$0{,}4$: wajar untuk pilihan ganda 4 opsi (tebakan acak 0,25).
      \item 81 soal sangat tajam (daya beda $>0{,}6$); 30 soal nyaris tidak berguna (daya beda $<0{,}1$).
      \item 8 soal menyentuh batas $g=0{,}7$ dan 18 soal batas $s=0{,}7$: soal ``aneh'' yang mungkin salah label di
            Q-matrix atau kuncinya membingungkan; layak ditinjau guru.
      \item Model otomatis memberi bobot kecil pada soal yang lemah (Bagian~\ref{sec:worked}), jadi soal buruk tidak
            merusak diagnosis, hanya tidak menyumbang informasi.
    \end{itemize}
  \end{minipage}
  \caption{Guess dan slip seluruh 948 soal. Garis putus-putus: daya beda 0,1 (merah) dan 0,6 (hijau).}
  \label{fig:gs}
\end{figure}

% ======================================================================================================
\section{Mengapa ``probabilistik'': dua lapis prior}

DINA klasik hanya mengestimasi $g$, $s$, dan proporsi penguasaan. P-DINA menambahkan dua \emph{prior}, yaitu
keyakinan awal yang dipakai sebelum melihat data.

\subsection{Prior untuk guess dan slip}

Setiap $g_j$ dan $s_j$ diberi prior $\mathrm{Beta}(1+0{,}2\kappa;\ 1+0{,}8\kappa)$. Cara mudah membacanya: prior ini
seperti menambahkan $\kappa$ ``jawaban bayangan'' dengan tingkat tebakan/kecerobohan 20\% ke setiap soal. Untuk soal
dengan ribuan jawaban, efeknya hampir nol; untuk soal dengan sedikit jawaban, prior mencegah nilai ekstrem (misalnya
$g=0$ hanya karena 3 siswa kebetulan salah semua). Model terpilih memakai $\kappa=2$, yaitu $\mathrm{Beta}(1{,}4;\
2{,}6)$. Selain itu ada batas keras: $g,s\in[0{,}001;\ 0{,}7]$ dan $g+s\le0{,}98$, agar soal selalu lebih mudah bagi
siswa yang menguasai.

\subsection{Prior higher-order untuk penguasaan}

Siswa yang kuat di banyak konsep cenderung kuat juga di konsep lain. P-DINA memodelkannya dengan satu angka per siswa,
\textbf{kemampuan umum} $\theta_i$, dan dua angka per konsep:
\[P(\alpha_{ik}=1\mid\theta_i)=\sigma(a_k\theta_i-b_k),\qquad \theta_i\sim\mathcal N(0,1),\qquad
\sigma(z)=\frac{1}{1+e^{-z}}.\]
$b_k$ adalah \emph{kesulitan} konsep (makin besar, makin sulit dikuasai) dan $a_k$ adalah seberapa kuat konsep itu
mengikuti kemampuan umum (makin besar, kurva makin curam).

\begin{figure}[h]
  \centering
  \setlength\figH{5.2cm}
  % Higher-order prior sigma(a_k*theta - b_k) for four target concepts (artifacts/model/pdina_params.npz)
% and the distribution of theta for the 3,934 train students (refit, slides/data/pdina_deep/theta_hist.csv).
\begin{tikzpicture}
\begin{axis}[
  name=curves, width=0.52\linewidth, height=\figH, xmin=-2.5, xmax=2.5, ymin=0, ymax=1,
  xlabel={kemampuan umum $\theta$}, ylabel={peluang awal menguasai},
  grid=major, legend style={at={(0.5,-0.28)}, anchor=north, font=\tiny, legend columns=2, /tikz/every even column/.append style={column sep=6pt}}, legend cell align=left,
  yticklabel style={/pgf/number format/fixed, /pgf/number format/precision=1},
]
  \addplot[navy, thick, domain=-2.5:2.5, samples=101] {1/(1+exp(-(5.000*x+0.665)))};
  \addlegendentry{Writing Expressions ($a{=}5{,}00$; $b{=}{-}0{,}67$)}
  \addplot[tealc, thick, domain=-2.5:2.5, samples=101] {1/(1+exp(-(4.666*x-0.978)))};
  \addlegendentry{Linear Equations ($4{,}67$; $0{,}98$)}
  \addplot[accent, thick, domain=-2.5:2.5, samples=101] {1/(1+exp(-(3.878*x-2.073)))};
  \addlegendentry{Horizontal \& Vertical Lines ($3{,}88$; $2{,}07$)}
  \addplot[gapred, thick, dashed, domain=-2.5:2.5, samples=101] {1/(1+exp(-(2.455*x-0.135)))};
  \addlegendentry{Gradient (2 soal; $2{,}46$; $0{,}14$)}
\end{axis}
\begin{axis}[
  at={(curves.east)}, anchor=west, xshift=1.0cm, ybar interval, width=0.34\linewidth, height=\figH,
  xmin=-2.5, xmax=2.5, ymin=0, xlabel={$\theta$ siswa train}, ylabel={jumlah siswa},
  ymajorgrids, xtick={-2,-1,...,2}, xticklabel={\pgfmathprintnumber{\tick}},
]
  \addplot[fill=navy!45, draw=white] coordinates {
    (-2.5,0) (-2.0,3) (-1.5,505) (-1.0,822) (-0.5,1078) (0.0,884) (0.5,465) (1.0,173) (1.5,4) (2.0,0) (2.5,0)};
\end{axis}
\end{tikzpicture}

  \caption{Kiri: peluang awal menguasai sebagai fungsi kemampuan umum untuk empat konsep target. Kanan: sebaran
  $\theta$ ke-3.934 siswa train.}
  \label{fig:ho}
\end{figure}

Gambar~\ref{fig:ho} memperlihatkan bahwa untuk siswa rata-rata ($\theta=0$) peluang awal menguasai \emph{Writing
Expressions} 0,66, sedangkan \emph{Horizontal and Vertical Lines} hanya 0,11: konsep kedua jauh lebih sulit. Manfaat
utamanya: konsep dengan sedikit jawaban ``meminjam'' informasi dari konsep lain milik siswa yang sama melalui
$\theta_i$, sehingga tidak hanya bertumpu pada 2--3 jawaban.

\begin{awas}
  Rata-rata $a_k$ untuk 57 konsep adalah 3,67, artinya penguasaan antar-konsep sangat terkait dengan kemampuan umum.
  Pada 4 dari 11 konsep target, $a_k$ menyentuh batas atas 5 yang dipasang di kode. Akibatnya kurva prior sangat
  curam: siswa dengan $\theta$ sedikit di bawah rata-rata langsung mendapat prior rendah. Ini salah satu alasan
  P(mastery) di EEDI cenderung sangat dekat ke 0 atau 1 (Bagian~\ref{sec:label}).
\end{awas}

% ======================================================================================================
\section{Dari jawaban ke P(mastery): menimbang bukti}
\label{sec:worked}

\subsection{Rumusnya ternyata sederhana}

Untuk soal yang menguji satu konsep (934 dari 948 soal), perhitungan P(mastery) adalah aturan Bayes biasa. Paling
mudah ditulis dalam \emph{log-odds} (logit), yaitu $\mathrm{logit}(p)=\ln\frac{p}{1-p}$:
\[
  \mathrm{logit}\,\pi_{ik} \;=\; \underbrace{\mathrm{logit}\,\sigma(a_k\theta_i-b_k)}_{\text{prior}}
  \;+\; \sum_{\text{jawaban benar}}\underbrace{\ln\frac{1-s_j}{g_j}}_{\text{bobot benar}}
  \;+\; \sum_{\text{jawaban salah}}\underbrace{\ln\frac{s_j}{1-g_j}}_{\text{bobot salah}} .
\]
Jadi setiap jawaban menambah atau mengurangi ``skor keyakinan'', dan \textbf{besarnya tergantung soalnya}.

\begin{table}[h]
  \centering\small
  \begin{tabularx}{\linewidth}{@{}lccrrL@{}}
    \toprule
    Soal & $g$ & $s$ & Bobot benar & Bobot salah & Artinya \\
    \midrule
    842 & 0,16 & 0,09 & $+1{,}73$ & $-2{,}18$ & soal tajam: setiap jawaban sangat informatif \\
    median 948 soal & 0,32 & 0,22 & $+0{,}77$ & $-1{,}14$ & soal tipikal \\
    820 & 0,26 & 0,55 & $+0{,}54$ & $-0{,}29$ & sering salah walau paham: jawaban salah lemah \\
    239 & 0,56 & 0,42 & $+0{,}04$ & $-0{,}05$ & hampir tidak membedakan apa pun \\
    \bottomrule
  \end{tabularx}
  \caption{Bobot bukti satu jawaban (dalam log-odds) untuk beberapa soal nyata. Baris median memakai median bobot
  dari 948 soal.}
  \label{tab:bobot}
\end{table}

Temuan penting dari Tabel~\ref{tab:bobot}: pada soal tipikal EEDI, \textbf{satu jawaban salah adalah bukti yang lebih
kuat ($-1{,}14$) daripada satu jawaban benar ($+0{,}77$)}. Alasannya, jawaban benar bisa saja hasil menebak (tebakan
sekitar 32\%), sedangkan siswa yang paham jarang salah. Karena itu 50\% benar \emph{tidak} berarti P(mastery) 0,5.

\subsection{Contoh nyata: satu siswa, satu konsep}

Siswa \texttt{Siswa-e7f6c011} (siswa test EEDI) menjawab 25 soal \emph{Writing Expressions}: 13 benar, 12 salah.
Kemampuan umumnya $\theta=-0{,}57$ (di bawah rata-rata, dari jawabannya pada konsep lain), sehingga prior-nya
$\sigma(5\times(-0{,}57)+0{,}67)=0{,}10$, atau $-2{,}19$ dalam log-odds.

\begin{figure}[h]
  \centering
  \setlength\figH{4.4cm}
  \begin{minipage}[c]{0.6\linewidth}% Worked example: log-odds contribution of each of the 25 Writing Expressions answers of Siswa-e7f6c011
% (slides/data/pdina_deep/worked_items.csv), sorted; plus the prior and the total.
\begin{tikzpicture}
\begin{axis}[
  ybar, width=\linewidth, height=\figH, bar width=5pt, ymin=-2.4, ymax=1.5,
  xmin=0.2, xmax=25.8, xtick=\empty, ylabel={sumbangan log-odds},
  ymajorgrids, xlabel={25 jawaban siswa, diurutkan (hijau = benar, merah = salah)},
  yticklabel style={/pgf/number format/fixed, /pgf/number format/precision=1},
]
  \addplot+[fill=okgreen, draw=none, bar shift=0pt] coordinates {
    (1,1.1669) (2,1.1390) (3,0.9155) (4,0.9150) (5,0.8992) (6,0.8770) (7,0.8518) (8,0.8063) (9,0.7962)
    (10,0.7447) (11,0.6078) (12,0.5905) (13,0.5074)};
  \addplot+[fill=gapred, draw=none, bar shift=0pt] coordinates {
    (14,-0.0691) (15,-0.1638) (16,-0.2854) (17,-0.5685) (18,-0.5733) (19,-0.5995) (20,-0.6837) (21,-0.7035)
    (22,-1.1678) (23,-1.2549) (24,-1.4057) (25,-1.8364)};
  \draw[black!60] (axis cs:0.2,0) -- (axis cs:25.8,0);
\end{axis}
\end{tikzpicture}
\end{minipage}\hfill
  \begin{minipage}[c]{0.37\linewidth}\resizebox{\linewidth}{!}{% Worked example as a waterfall: prior logit + evidence of correct + evidence of wrong = posterior logit.
\begin{tikzpicture}[x=1.35cm, y=0.2cm, font=\scriptsize]
  \foreach \v in {-10,-5,0,5,10} {\draw[midgray!40] (-0.3,\v) -- (4.3,\v); \node[anchor=east] at (-0.35,\v) {$\v$};}
  \draw[midgray] (-0.3,0) -- (4.3,0);
  \node[rotate=90, anchor=south] at (-1.0,0) {log-odds};
  \fill[midgray] (0,0) rectangle (0.7,-2.19);
  \fill[okgreen] (1,-2.19) rectangle (1.7,8.63);
  \fill[gapred] (2,8.63) rectangle (2.7,-0.68);
  \fill[navy] (3.3,0) rectangle (4.0,-0.68);
  \draw[dashed, black!50] (0.7,-2.19) -- (1,-2.19) (1.7,8.63) -- (2,8.63) (2.7,-0.68) -- (3.3,-0.68);
  \foreach \x/\t in {0.35/{prior\\$-2{,}19$}, 1.35/{13 benar\\$+10{,}82$}, 2.35/{12 salah\\$-9{,}31$}, 3.65/{akhir\\$-0{,}68$}}
    \node[align=center, anchor=north] at (\x,-10.8) {\t};
  \node[anchor=south, text=navy] at (3.65,0.4) {$\pi=0{,}34$};
\end{tikzpicture}
}\end{minipage}
  \caption{Kiri: bobot setiap jawaban siswa. Kanan: prior ditambah total bukti benar dan salah menghasilkan log-odds
  akhir $-0{,}68$, yaitu P(mastery) = 0,34.}
  \label{fig:worked}
\end{figure}

Hitungan lengkapnya: $-2{,}19+10{,}82-9{,}31=-0{,}68$, dan $\sigma(-0{,}68)=0{,}3357$. Fungsi \texttt{infer()} model
memberi \textbf{0,3357} juga; rumus di atas memang persis yang dikerjakan model untuk soal satu-konsep. Perhatikan:
13 jawaban benar menyumbang $+10{,}82$, tetapi 12 jawaban salah menyumbang $-9{,}31$ dan prior awalnya sudah negatif.
Hasil akhirnya, siswa ini lebih mungkin \emph{belum} menguasai, walaupun lebih dari separuh jawabannya benar.

\subsection{Seberapa cepat keyakinan berubah}

Gambar~\ref{fig:traj} memperlihatkan siswa baru (tanpa jawaban di konsep lain, jadi $\theta$ hanya ditaksir dari
jawaban ini) yang mengerjakan 10 soal \emph{Writing Expressions} dengan daya beda tipikal.

\begin{figure}[h]
  \centering
  \setlength\figH{5cm}
  \begin{minipage}{0.75\linewidth}% P(mastery) of a new student after n answers on 10 typical Writing Expressions items
% (slides/data/pdina_deep/trajectory.csv; computed with the trained model's infer()).
\begin{tikzpicture}
\begin{axis}[
  width=\linewidth, height=\figH, xmin=-0.3, xmax=10.3, ymin=-0.03, ymax=1.03,
  xlabel={jumlah jawaban pada konsep ini}, ylabel={P(mastery)}, grid=major, xtick={0,...,10},
  legend style={at={(0.98,0.82)}, anchor=north east, fill=white, fill opacity=0.9, text opacity=1}, legend cell align=left,
  yticklabel style={/pgf/number format/fixed, /pgf/number format/precision=1},
]
  \fill[accent!10] (axis cs:-0.3,-0.03) rectangle (axis cs:4.5,1.03);
  \node[font=\tiny, text=accent!80!black] at (axis cs:2.1,0.2) {bukti $<n_{\min}=5$};
  \addplot[okgreen, thick, mark=*, mark size=1.3pt] table[col sep=comma, x=n, y=semua_benar] {data/pdina_deep/trajectory.csv};
  \addlegendentry{semua benar}
  \addplot[accent, thick, mark=diamond*, mark size=1.6pt] table[col sep=comma, x=n, y=selang_seling] {data/pdina_deep/trajectory.csv};
  \addlegendentry{benar--salah bergantian}
  \addplot[gapred, thick, mark=square*, mark size=1.2pt] table[col sep=comma, x=n, y=semua_salah] {data/pdina_deep/trajectory.csv};
  \addlegendentry{semua salah}
  \draw[dashed, black!60] (axis cs:-0.3,0.5) -- (axis cs:10.3,0.5);
  \node[font=\tiny, anchor=north east] at (axis cs:10.2,0.5) {$\tau=0{,}5$};
\end{axis}
\end{tikzpicture}
\end{minipage}
  \caption{P(mastery) setelah $n$ jawaban, dihitung dengan model terlatih.}
  \label{fig:traj}
\end{figure}

Tiga hal terlihat. Pertama, keyakinan bergerak sangat cepat: satu jawaban benar langsung menaikkan P(mastery) dari
0,66 ke 0,90. Kedua, pola benar--salah bergantian berakhir \emph{rendah} (0,003 setelah 10 jawaban), karena 50\% benar
pada soal dengan tebakan sekitar 40\% lebih mirip siswa yang menebak daripada siswa yang paham (yang benar sekitar
85\%). Ketiga, pada 1--4 jawaban pertama nilainya masih melompat-lompat. Itulah alasan ada aturan \textbf{bukti minimum}:
konsep dengan kurang dari 5 jawaban tidak diberi label gap, seberapa pun rendah probabilitasnya.

\subsection{Persentase benar bukan P(mastery)}

Apakah P-DINA sekadar ``persentase benar yang dihaluskan''? Gambar~\ref{fig:pcpi} membandingkan keduanya pada 2.121 sel
siswa test $\times$ konsep target dengan bukti minimal 5 jawaban.

\begin{figure}[h]
  \centering
  \setlength\figH{5.4cm}
  \begin{minipage}[c]{0.55\linewidth}% Percent correct vs P(mastery) for 2,121 test cells (target concepts, >= 5 answers);
% slides/data/pdina_deep/pc_vs_pi_test.csv, from artifacts/diagnosis/test_mastery_profiles.csv
\begin{tikzpicture}
\begin{axis}[
  width=\linewidth, height=\figH, xmin=-0.02, xmax=1.02, ymin=-0.03, ymax=1.03,
  xlabel={persentase jawaban benar pada konsep}, ylabel={P(mastery) dari P-DINA}, grid=major,
  xtick={0,0.2,...,1}, ytick={0,0.25,0.5,0.75,1},
  xticklabel={\pgfmathparse{\tick*100}\pgfmathprintnumber[fixed,precision=0]{\pgfmathresult}\%},
  yticklabel style={/pgf/number format/fixed, /pgf/number format/precision=2},
]
  \fill[accent!12] (axis cs:0.5,-0.03) rectangle (axis cs:0.6,1.03);
  \addplot[only marks, mark=*, mark size=0.8pt, navy, fill opacity=0.3, draw opacity=0.3]
    table[col sep=comma, x=pc, y=mastery_prob] {data/pdina_deep/pc_vs_pi_test.csv};
  \draw[dashed, gapred] (axis cs:-0.02,0.5) -- (axis cs:1.02,0.5);
  \node[font=\tiny, text=gapred, anchor=south west] at (axis cs:0.0,0.5) {$\tau=0{,}5$};
  \node[font=\tiny, text=accent!80!black, anchor=north, align=center, fill=white, inner sep=1pt] at (axis cs:0.55,0.42)
    {50--60\% benar:\\69\% tetap gap};
\end{axis}
\end{tikzpicture}
\end{minipage}\hfill
  \begin{minipage}[c]{0.42\linewidth}\small
    \begin{tabular}{@{}lrr@{}}
      \toprule
      \% benar & sel & label gap \\
      \midrule
      40--50\% & 232 & 94\% \\
      50--60\% & 274 & 69\% \\
      60--70\% & 211 & 28\% \\
      70--80\% & 128 & 3\% \\
      \bottomrule
    \end{tabular}\\[4pt]
    Korelasi keduanya 0,85: searah, tetapi tidak sama. Batas ``paham'' versi P-DINA berada di sekitar 60\%
    benar, bukan 50\%, karena tebakan membuat 30--40\% benar bisa dicapai tanpa paham.
  \end{minipage}
  \caption{Persentase benar vs P(mastery) pada siswa test (\texttt{test\_mastery\_profiles.csv}).}
  \label{fig:pcpi}
\end{figure}

Sebaran vertikal yang lebar pada persentase yang sama terjadi karena dua hal: soal yang dikerjakan berbeda (benar di
soal tajam bernilai lebih dari benar di soal yang mudah ditebak), dan kemampuan umum $\theta$ siswa ikut
memengaruhi prior. Inilah nilai tambah P-DINA dibanding skor proporsi: dua siswa dengan 55\% benar bisa didiagnosis
berbeda, dan alasannya dapat ditelusuri.

\textbf{Soal multi-konsep.} Untuk 14 soal yang menguji dua konsep, siswa siap hanya jika menguasai keduanya, sehingga
peluang siap $\rho_{ij}=\prod_k\pi_{ik}^{q_{jk}}$. Bukti dari soal itu dibagi ke kedua konsep secara proporsional, dan
karena bagiannya saling bergantung, dihitung berulang (iterasi titik tetap dengan \emph{damping} 0,5) sampai stabil.

% ======================================================================================================
\section{Training: variational EM}

\subsection{Mengapa tidak dihitung langsung}

Bagian sebelumnya mengandaikan $g$, $s$, $a$, $b$ sudah diketahui. Saat training, keduanya harus dicari bersamaan:
parameter soal butuh status penguasaan siswa, dan status penguasaan butuh parameter soal. Algoritme EM menyelesaikan
``ayam dan telur'' ini dengan bergantian menebak satu sisi lalu memperbaiki sisi lain.

EM klasik untuk DINA harus menghitung peluang setiap kombinasi penguasaan. Dengan 57 konsep ada
$2^{57}\approx1{,}4\times10^{17}$ kombinasi per siswa, mustahil. P-DINA memakai pendekatan \textbf{mean-field}: status
setiap konsep dianggap berdiri sendiri, $q(\boldsymbol\alpha_i)=\prod_k\mathrm{Bern}(\pi_{ik})$, sehingga cukup
menyimpan 57 angka per siswa (Gambar~\ref{fig:mf}). Untuk soal satu-konsep pendekatan ini tidak kehilangan apa pun;
aproksimasi hanya terjadi pada 14 soal multi-konsep.

\begin{figure}[h]
  \centering
  \resizebox{0.92\linewidth}{!}{% Illustration: exact posterior over 2^K profiles vs. mean-field (K independent Bernoullis), drawn for K = 3.
\begin{tikzpicture}[cell/.style={draw=midgray, minimum width=7mm, minimum height=5mm, font=\tiny\ttfamily, inner sep=0pt}]
  \node[font=\scriptsize\bfseries, text=navy, anchor=west] at (0,1.0) {Posterior eksak: satu peluang per profil};
  \foreach \p [count=\i from 0] in {000,001,010,011,100,101,110,111} {
    \node[cell, fill=navy!\fpeval{8+\i*9}] at (\i*0.8+0.35,0.3) {\p};
  }
  \node[font=\tiny, text=black!70, anchor=west] at (0,-0.35) {$2^3=8$ angka; untuk 57 konsep: $2^{57}\approx1{,}4\times10^{17}$ angka per siswa};

  \draw[arr] (6.9,0.3) -- node[lbl, above]{mean-field} (8.3,0.3);

  \node[font=\scriptsize\bfseries, text=accent!80!black, anchor=west] at (8.5,1.0) {Mean-field: satu peluang per konsep};
  \foreach \k/\v [count=\i from 0] in {A/0.85,B/0.10,C/0.60} {
    \node[font=\tiny, anchor=east] at (9.0,0.45-\i*0.38) {\k};
    \draw[midgray, fill=lightbg] (9.1,0.35-\i*0.38) rectangle (11.6,0.55-\i*0.38);
    \fill[accent] (9.1,0.35-\i*0.38) rectangle ({9.1+2.5*\v},0.55-\i*0.38);
    \node[font=\tiny, anchor=west] at (11.65,0.45-\i*0.38) {$\pi=\v$};
  }
  \node[font=\tiny, text=black!70, anchor=west] at (8.5,-0.75) {3 angka; untuk 57 konsep: 57 angka per siswa};
\end{tikzpicture}
}
  \caption{Mean-field mengganti tabel peluang $2^K$ profil dengan $K$ peluang terpisah.}
  \label{fig:mf}
\end{figure}

Yang dioptimalkan adalah \textbf{ELBO} (\emph{evidence lower bound}), batas bawah log-likelihood data yang punya
bentuk tertutup di bawah mean-field:
\[
  \mathrm{ELBO}=\sum_{i,j} m_{ij}\big[\rho_{ij}\,\ell^{(1)}_{ij}+(1-\rho_{ij})\,\ell^{(0)}_{ij}\big]
  -\sum_{i,k}\mathrm{KL}\big(\mathrm{Bern}(\pi_{ik})\,\|\,\mathrm{Bern}(\lambda_{ik})\big)+\log p(g,s,\theta),
\]
dengan $m_{ij}=1$ bila jawaban teramati, $\ell^{(1)}$ dan $\ell^{(0)}$ log-likelihood jawaban bila siswa siap atau
belum siap, dan $\lambda_{ik}=\sigma(a_k\theta_i-b_k)$. Suku pertama: ``seberapa baik model menjelaskan jawaban''.
Suku kedua: ``jangan terlalu jauh dari prior tanpa bukti''. \textbf{Hanya jawaban teramati yang dijumlahkan}, jadi
soal kosong tidak pernah dihitung sebagai salah.

\subsection{Satu putaran EM}

\begin{figure}[h]
  \centering
  \resizebox{0.95\linewidth}{!}{% Variational EM training loop.
\begin{tikzpicture}[node distance=4mm and 5mm]
  \node[boxsoft, text width=2.3cm] (init) {\textbf{Inisialisasi}\\{\tiny $g_j{=}s_j{=}0{,}2$, $a_k{=}1$, $b_k{=}0$, $\theta_i{=}0$}};
  \node[boxnavy, right=of init, text width=2.5cm] (e) {\textbf{E-step}\\{\tiny perbarui $\pi_{ik}$ (keyakinan mastery) dan $\theta_i$ (Newton)}};
  \node[boxnavy, right=of e, text width=2.5cm] (m) {\textbf{M-step}\\{\tiny $g_j, s_j$ tertutup (MAP Beta); $a_k, b_k$ (Newton)}};
  \node[boxsoft, right=of m, text width=2.1cm] (elbo) {\textbf{Hitung ELBO}};
  \node[diamond, draw=navy, aspect=2.2, inner sep=1pt, font=\tiny, align=center, right=of elbo] (conv) {$\Delta\text{ELBO}/|\text{ELBO}|$\\$\le 10^{-6}$?};
  \node[boxacc, below=7mm of conv, text width=1.9cm] (stop) {\textbf{Checkpoint}\\{\tiny $g, s, a, b$}};
  \draw[arr] (init) -- (e);
  \draw[arr] (e) -- (m);
  \draw[arr] (m) -- (elbo);
  \draw[arr] (elbo) -- (conv);
  \draw[arr] (conv) -- node[lbl, right]{ya (atau 150 iterasi)} (stop);
  \draw[arr] (conv.north) -- ++(0,0.55) -| node[lbl, pos=0.25, above]{tidak} (e.north);
\end{tikzpicture}
}
  \caption{Alur training. Semua langkah memakai perkalian matriks, tanpa GPU.}
\end{figure}

\begin{itemize}
  \item \textbf{Mulai}: $g_j=s_j=0{,}2$, $a_k=1$, $b_k=0$, $\theta_i=0$.
  \item \textbf{E-step}: dengan $g,s,a,b$ tetap, hitung $\pi_{ik}$ semua siswa memakai rumus log-odds di Bagian~\ref{sec:worked}
        (dua perkalian matriks untuk seluruh siswa sekaligus), lalu perbarui $\theta_i$ dengan metode Newton.
  \item \textbf{M-step}: dengan $\pi$ tetap, perbarui parameter soal. Rumusnya tertutup dan mudah dibaca:
        \[
          s_j=\frac{\text{(siswa siap yang salah)}+a_s-1}{\text{(siswa siap)}+a_s+b_s-2},\qquad
          g_j=\frac{\text{(siswa belum siap yang benar)}+a_g-1}{\text{(siswa belum siap)}+a_g+b_g-2},
        \]
        di mana ``siswa siap'' dihitung secara lunak (dijumlahkan dengan bobot $\rho_{ij}$). Jadi slip adalah
        proporsi salah di antara yang (diperkirakan) siap, ditambah jawaban bayangan dari prior. Lalu $a_k,b_k$
        diperbarui dengan Newton.
  \item \textbf{Berhenti} bila ELBO naik kurang dari $10^{-6}$ relatif, atau setelah 150 putaran.
\end{itemize}

\begin{figure}[h]
  \centering
  \setlength\figH{4.3cm}
  \begin{minipage}{0.7\linewidth}% ELBO per EM iteration, refit of the selected configuration on train (slides/data/pdina_deep/elbo_history.csv)
\begin{tikzpicture}
\begin{axis}[
  width=\linewidth, height=\figH, xmin=0, xmax=45, xlabel={iterasi EM}, ylabel={ELBO (ribuan)},
  grid=major, yticklabel style={/pgf/number format/fixed, /pgf/number format/precision=0},
]
  \addplot[navy, thick, mark=*, mark size=0.9pt] table[col sep=comma, x=iter, y expr=\thisrow{elbo}/1000] {data/pdina_deep/elbo_history.csv};
  \node[font=\tiny, anchor=south east, align=right] at (axis cs:44,-643.5) {iterasi 44: $-630.546$\\perubahan relatif $\le10^{-6}$ $\Rightarrow$ berhenti};
\end{axis}
\end{tikzpicture}
\end{minipage}
  \caption{ELBO per iterasi pada konfigurasi terpilih. Kenaikan terbesar terjadi di 10 iterasi pertama.}
\end{figure}

Model terpilih berhenti pada iterasi 44 dengan ELBO $-630.546$, dalam 32 detik di CPU laptop. Karena tidak ada unsur
acak dalam training, menjalankan ulang memberi parameter yang persis sama dengan checkpoint.

% ======================================================================================================
\section{Dari probabilitas ke label}
\label{sec:label}

Guru butuh keputusan, bukan hanya angka. P(mastery) diubah menjadi label dengan tiga angka kebijakan yang dipilih dari
data validation (Gambar~\ref{fig:policy}).

\begin{figure}[h]
  \centering
  \resizebox{0.98\linewidth}{!}{% Labelling rule: evidence first, then probability.
\begin{tikzpicture}[node distance=5mm and 9mm]
  \node[boxsoft, text width=2.3cm] (in) {P(mastery) $\pi$\\ dan jumlah bukti $n$};
  \node[diamond, draw=navy, aspect=2, inner sep=1pt, font=\tiny, align=center, right=of in] (q1) {$n\ge n_{\min}$?\\($n_{\min}=5$)};
  \node[box, draw=midgray, below=8mm of q1, text width=2.6cm] (ins) {\textbf{insufficient\_evidence}\\{\tiny butuh soal tambahan; \emph{bukan} gap}};
  \node[diamond, draw=navy, aspect=2, inner sep=1pt, font=\tiny, align=center, right=of q1] (q2) {$|\pi-\tau|<\delta$?\\($\delta=0$)};
  \node[box, draw=accent, below=8mm of q2, text width=2.2cm] (unc) {\textbf{uncertain}\\{\tiny terlalu dekat batas}};
  \node[diamond, draw=navy, aspect=2, inner sep=1pt, font=\tiny, align=center, right=of q2] (q3) {$\pi\ge\tau$?\\($\tau=0{,}5$)};
  \node[box, draw=okgreen, above right=2mm and 6mm of q3, text width=1.9cm] (mas) {\textbf{mastered}};
  \node[box, draw=gapred, below right=2mm and 6mm of q3, text width=1.9cm] (gap) {\textbf{not\_mastered}\\{\tiny = learning gap}};
  \draw[arr] (in) -- (q1);
  \draw[arr] (q1) -- node[lbl, right]{tidak} (ins);
  \draw[arr] (q1) -- node[lbl, above]{ya} (q2);
  \draw[arr] (q2) -- node[lbl, right]{ya} (unc);
  \draw[arr] (q2) -- node[lbl, above]{tidak} (q3);
  \draw[arr] (q3) |- node[lbl, pos=0.25, left]{ya} (mas);
  \draw[arr] (q3) |- node[lbl, pos=0.25, left]{tidak} (gap);
\end{tikzpicture}
}
  \caption{Urutan aturan label. Jumlah bukti dicek lebih dulu.}
  \label{fig:policy}
\end{figure}

\begin{itemize}
  \item $\boldsymbol{n_{\min}=5}$: respons tiap siswa dibagi dua secara acak, diagnosis dibuat dari masing-masing
        paruh, lalu dihitung seberapa sering keduanya sepakat. Lima adalah jumlah terkecil yang mencapai kesepakatan
        minimal 0,85.
  \item $\boldsymbol{\tau=0{,}5}$: semua nilai $\tau$ antara 0,45 dan 0,70 memberi prediksi yang setara (dalam satu
        standard error dari yang terbaik), sehingga dipilih nilai netral 0,5.
  \item $\boldsymbol{\delta=0}$: tanpa pita ketidakpastian target kesepakatan sudah tercapai. Gambar~\ref{fig:pihist}
        menjelaskan alasannya: P(mastery) hampir selalu dekat 0 atau 1.
\end{itemize}

\begin{figure}[h]
  \centering
  \setlength\figH{4.2cm}
  \begin{minipage}[c]{0.58\linewidth}% Distribution of P(mastery) on train, target concepts with >= 5 responses (16,410 cells; pi_hist_target.csv)
\begin{tikzpicture}
\begin{axis}[
  ybar interval, width=\linewidth, height=\figH, xmin=0, xmax=1, ymode=log, ymin=100, ymax=15000, log origin=infty,
  xlabel={P(mastery)}, ylabel={jumlah sel siswa $\times$ konsep}, ymajorgrids,
  xtick={0,0.1,...,1}, xticklabel={\pgfmathprintnumber[fixed,precision=1]{\tick}},
]
  \addplot[fill=navy!55, draw=white] coordinates {
    (0.0,9290) (0.1,374) (0.2,244) (0.3,221) (0.4,160) (0.5,196) (0.6,197) (0.7,209) (0.8,375) (0.9,5144) (1.0,100)};
\end{axis}
\end{tikzpicture}
\end{minipage}\hfill
  \begin{minipage}[c]{0.38\linewidth}\small
    Dari 16.410 sel siswa train $\times$ konsep target dengan bukti $\ge5$:\\[3pt]
    \begin{tabular}{@{}lr@{}}
      \toprule
      P(mastery) $<0{,}1$ & 56,6\% \\
      P(mastery) $>0{,}9$ & 31,3\% \\
      $0{,}4<$ P(mastery) $<0{,}6$ & 2,2\% \\
      \bottomrule
    \end{tabular}
  \end{minipage}
  \caption{Sebaran P(mastery) di train sangat terpolarisasi, sehingga pilihan $\tau$ dan $\delta$ hanya memengaruhi
  sedikit sel.}
  \label{fig:pihist}
\end{figure}

\begin{awas}
  \textbf{Membaca kolom SD.} Output menyertakan $\mathrm{SD}=\sqrt{\pi(1-\pi)}$. Angka ini sepenuhnya ditentukan oleh
  $\pi$ (maksimum 0,5 saat $\pi=0{,}5$) dan \emph{tidak} memperhitungkan jumlah bukti maupun ketidakpastian parameter
  soal. Konsep dengan 4 jawaban dan $\pi=0{,}00$ tetap punya SD sekitar 0. Untuk menilai seberapa kuat dasar
  diagnosis, lihat kolom jumlah bukti, bukan SD.
\end{awas}

% ======================================================================================================
\section{Seberapa bagus modelnya?}

EEDI tidak punya ``kunci jawaban'' penguasaan, jadi kualitas diukur secara tidak langsung, semuanya pada siswa yang
tidak dipakai untuk training.

\begin{table}[h]
  \centering\small
  \begin{tabularx}{\linewidth}{@{}L r L@{}}
    \toprule
    Ukuran (test, 492 siswa) & Hasil & Cara membaca \\
    \midrule
    ROC-AUC prediksi 28.003 jawaban tersembunyi & 0,768 & setara Q-logistic (0,775); log-loss tidak berbeda bermakna \\
    ECE (kalibrasi) & 0,008 & ``P(benar) 0,7'' memang benar sekitar 70\%; Q-logistic 0,036 \\
    Konsistensi split-half (880 sel) & 0,876 & dua diagnosis dari data berbeda sepakat 88\% \\
    Akurasi atribut, simulasi (ada / tanpa bukti) & 0,957 / 0,851 & batas atas: data dibangkitkan dari model sendiri \\
    Recovery parameter, simulasi $r(g)$ / $r(s)$ & 0,965 / 0,982 & prosedur training menemukan kembali parameter yang benar \\
    \bottomrule
  \end{tabularx}
  \caption{Ringkasan evaluasi. Detail dan grafik ada di \texttt{README\_pdina.md} dan slide presentasi.}
\end{table}

Prediksi jawaban yang baik tidak membuktikan diagnosis benar, karena itu ketiga jenis ukuran dilaporkan bersama.

% ======================================================================================================
\section{Memakai model: siswa baru dan kapan training ulang}

Untuk siswa baru, model \textbf{tidak dilatih ulang}. Fungsi \texttt{infer()} hanya menjalankan E-step dengan
$g,s,a,b$ yang dibekukan, sehingga diagnosis 30 siswa butuh sekitar 0,4 detik dan bisa memakai tes yang tidak lengkap.
Syaratnya, soal yang dipakai harus sudah ada di item bank checkpoint.

\begin{table}[h]
  \centering\small
  \begin{tabularx}{\linewidth}{@{}L l L@{}}
    \toprule
    Situasi & Training ulang? & Alasan \\
    \midrule
    Siswa baru mengerjakan soal yang sudah ada & tidak & cukup E-step \\
    Mengubah $\tau$, $\delta$, $n_{\min}$ atau peta intervensi & tidak & hanya kebijakan, bukan model \\
    Menambah soal baru & ya & $g,s$ soal baru harus dikalibrasi (sekitar $\ge30$ jawaban per soal) \\
    Mengubah Q-matrix atau daftar konsep & ya & struktur model berubah \\
    Populasi lain (kelas SMA penelitian) & ya / validasi ulang & parameter EEDI berasal dari siswa SMP di Inggris \\
    \bottomrule
  \end{tabularx}
  \caption{Kapan checkpoint perlu dilatih ulang.}
\end{table}

% ======================================================================================================
\section{Asumsi dan keterbatasan}

\begin{table}[h]
  \centering\small
  \begin{tabularx}{\linewidth}{@{}l L L@{}}
    \toprule
    Asumsi & Artinya & Akibat bila tidak terpenuhi \\
    \midrule
    Q-matrix benar & soal benar-benar menguji konsep yang tertulis & diagnosis salah secara sistematis; audit 22 soal: 17 cocok, 5 parsial \\
    Konjungtif (``AND'') & soal multi-konsep butuh semua konsepnya & terlalu ketat bila sebenarnya bisa saling menutupi (hanya 14 soal) \\
    Penguasaan statis & siswa tidak belajar selama tes & perubahan antar-waktu dianalisis terpisah (pre/post) \\
    Independensi lokal & bila status diketahui, jawaban saling bebas & soal yang mirip bisa ``dihitung dua kali'' \\
    Mean-field & status antar-konsep dihitung terpisah & eksak untuk soal satu-konsep; aproksimasi untuk 14 soal \\
    $\theta$ titik (MAP) & kemampuan umum tidak diberi rentang ketidakpastian & keyakinan bisa terlalu tinggi; diperkuat oleh $a_k$ yang menyentuh batas 5 \\
    \bottomrule
  \end{tabularx}
  \caption{Asumsi P-DINA dan dampaknya.}
\end{table}

Keterbatasan data: EEDI tidak punya ground truth penguasaan; populasinya siswa SMP di Inggris dengan 11 konsep target
yang hanya merupakan prasyarat aljabar SMA (tanpa kalkulus); daya beda soal sedang sehingga banyak sel masih
\emph{insufficient\_evidence}; dan \emph{Gradient Between Two Co-ordinates} hanya punya 2 soal.

\begin{intinya}
  \textbf{Ringkasnya.} P-DINA menimbang setiap jawaban sesuai seberapa informatif soalnya (guess dan slip),
  menggabungkannya dengan kemampuan umum siswa, dan menghasilkan P(mastery) yang terkalibrasi dengan sekitar 2.000
  parameter yang semuanya bisa dijelaskan. Kekuatannya adalah transparansi: setiap angka dapat ditelusuri sampai ke
  jawaban yang membentuknya, seperti contoh 0,34 di Bagian~\ref{sec:worked}.
\end{intinya}

% ======================================================================================================
\appendix
\section{Daftar notasi}

\begin{table}[h]
  \centering\small
  \begin{tabularx}{\linewidth}{@{}l L l L@{}}
    \toprule
    Simbol & Arti & Simbol & Arti \\
    \midrule
    $X_{ij}$ & jawaban siswa $i$ pada soal $j$ (1/0) & $m_{ij}$ & 1 bila jawaban teramati \\
    $q_{jk}$ & 1 bila soal $j$ menguji konsep $k$ & $\alpha_{ik}$ & penguasaan sebenarnya (tidak teramati) \\
    $\eta_{ij}$ & siswa siap untuk soal (semua konsep dikuasai) & $\pi_{ik}$ & P(mastery), output model \\
    $\rho_{ij}$ & peluang siswa siap, $\prod_k\pi_{ik}^{q_{jk}}$ & $\lambda_{ik}$ & peluang awal, $\sigma(a_k\theta_i-b_k)$ \\
    $g_j$, $s_j$ & guess dan slip soal $j$ & $\theta_i$ & kemampuan umum siswa \\
    $a_k$, $b_k$ & keterkaitan dan kesulitan konsep $k$ & $\kappa$ & kekuatan prior Beta untuk $g,s$ \\
    $\tau$, $\delta$, $n_{\min}$ & threshold, pita ragu, bukti minimum & $\sigma(\cdot)$ & fungsi logistik \\
    \bottomrule
  \end{tabularx}
\end{table}

\sumber{Referensi: Junker \& Sijtsma (2001) untuk DINA; de la Torre \& Douglas (2004) untuk prior higher-order; Blei,
Kucukelbir \& McAuliffe (2017) untuk inferensi variasional; Wang et al.\ (2020) untuk data EEDI. Kode:
\texttt{cdacdm/models.py}, \texttt{cdacdm/diagnosis.py}; angka tambahan: \texttt{slides/data/pdina\_deep/}.}

\end{document}
